% GENERATED FILE. Edit canonical lessons, then run npm run generate:books.
% canonical-source-hash: fnv1a64:bde3d59a78cd00d2
% canonical-lessons: PA-C242-talna, PA-C242-bhunnna, PA-C242-karchhi, PA-C242-chhirakna, PA-C242-sarna

\chapter{Frying and Roasting}
\label{ch:pa-a2-frying-and-roasting}
\hlchaptermodality{242}

\begin{chapteropening}
\textbf{By the end of this chapter:} \emph{I can talk about frying and roasting food.}

Complete the last lesson of chapter 242: I can talk about frying and roasting food.
\end{chapteropening}

\section[\texorpdfstring{\pa{ਤਲਣਾ} (talṇā)}{ਤਲਣਾ (talṇā)}]{\pa{ਤਲਣਾ} (talṇā) — to fry}

\label{lesson:PA-C242-talna}



\begin{quote}
Before the new one: say the Punjabi for a chilli, then the Punjabi for gram flour.
\end{quote}

\subsection*{You'll want to know: \pa{ਤਲਣਾ}}
\textbf{\pa{ਤਲਣਾ}} — \emph{talṇā} — \textquotedblleft{}to fry\textquotedblright{}.

\begin{quote}
\pa{ਮੈਂ} \pa{ਪਕੌੜੇ} \pa{ਤੇਲ} \pa{ਵਿੱਚ} \pa{ਤਲਦੀ} \pa{ਹਾਂ।}
\end{quote}

\begin{quote}
\emph{maiṁ pakauṛe tel vich taldī hāṁ.}
\end{quote}

\begin{quote}
I fry the pakoras in oil.
\end{quote}

Past (I …): \pa{ਮੈਂ} \pa{ਪਕੌੜੇ} \pa{ਤਲੇ}, \emph{maiṁ pakauṛe tale}.

\begin{cousinweb}
\textbf{\pa{ਤਲਣਾ}} (\emph{talṇā}) is inherited from a reconstructed \emph{*talati}, \textquotedblleft{}fries\textquotedblright{}; CDIAL lists Punjabi \emph{talṇā}, \textquotedblleft{}to fry, cook in ghee\textquotedblright{} (CDIAL 5736).
\end{cousinweb}

\begin{grammarlens}[title={What you've built}]
One more word for this chapter.
\end{grammarlens}

\subsection*{Your turn}
\begin{itemize}
\raggedright
  \item \emph{Say it:} \emph{talṇā}
  \item \emph{Say it:} \emph{talṇā}, once more
  \item \emph{Say it:} \emph{maiṁ pakauṛe tel vich taldī hāṁ.} — I fry the pakoras in oil.
  \item \emph{Recall:} say the Punjabi for a chilli, then the Punjabi for gram flour, then say \emph{talṇā} again
\end{itemize}

\begin{tcolorbox}[breakable,colback=teal!4,colframe=teal!35!black,title={Before you move on}]
What does \pa{ਤਲਣਾ} mean? (\textquotedblleft{}To fry\textquotedblright{}.) Say it once more.
\end{tcolorbox}
\section[\texorpdfstring{\pa{ਭੁੰਨਣਾ} (bhunnṇā)}{ਭੁੰਨਣਾ (bhunnṇā)}]{\pa{ਭੁੰਨਣਾ} (bhunnṇā) — to roast; to parch}

\label{lesson:PA-C242-bhunnna}



\begin{quote}
Before the new one: say the Punjabi for gram flour, then the Punjabi for to fry.
\end{quote}

\subsection*{You'll want to know: \pa{ਭੁੰਨਣਾ}}
\textbf{\pa{ਭੁੰਨਣਾ}} — \emph{bhunnṇā} — \textquotedblleft{}to roast; to parch\textquotedblright{}.

\begin{quote}
\pa{ਸਰਦੀ} \pa{ਵਿੱਚ} \pa{ਅਸੀਂ} \pa{ਮੂੰਗਫਲੀ} \pa{ਭੁੰਨਦੇ} \pa{ਹਾਂ।}
\end{quote}

\begin{quote}
\emph{sardī vich asīṁ mū̃gphalī bhunnde hāṁ.}
\end{quote}

\begin{quote}
In winter we roast peanuts.
\end{quote}

Past (I …): \pa{ਮੈਂ} \pa{ਮੱਕੀ} \pa{ਭੁੰਨੀ}, \emph{maiṁ makkī bhunnī}.

\begin{cousinweb}
\textbf{\pa{ਭੁੰਨਣਾ}} (\emph{bhunnṇā}) is inherited from a reconstructed \emph{*bhṛgṇa}, \textquotedblleft{}fried, roasted\textquotedblright{}; CDIAL lists Punjabi \emph{bhunnaṇā}, \textquotedblleft{}to parch, fry\textquotedblright{} (CDIAL 9577).
\end{cousinweb}

\begin{grammarlens}[title={What you've built}]
One more word for this chapter.
\end{grammarlens}

\subsection*{Your turn}
\begin{itemize}
\raggedright
  \item \emph{Say it:} \emph{bhunnṇā}
  \item \emph{Say it:} \emph{bhunnṇā}, once more
  \item \emph{Say it:} \emph{sardī vich asīṁ mū̃gphalī bhunnde hāṁ.} — In winter we roast peanuts.
  \item \emph{Recall:} say the Punjabi for gram flour, then the Punjabi for to fry, then say \emph{bhunnṇā} again
  \item \emph{Write it:} copy \textbf{\pa{ਭੁੰਨਣਾ}} once, checking the whole word
\end{itemize}

\begin{tcolorbox}[breakable,colback=teal!4,colframe=teal!35!black,title={Before you move on}]
What does \pa{ਭੁੰਨਣਾ} mean? (\textquotedblleft{}To roast; to parch\textquotedblright{}.) Say it once more.
\end{tcolorbox}
\section[\texorpdfstring{\pa{ਕੜਛੀ} (kaṛchhī)}{ਕੜਛੀ (kaṛchhī)}]{\pa{ਕੜਛੀ} (kaṛchhī) — a ladle}

\label{lesson:PA-C242-karchhi}



\begin{quote}
Before the new one: say the Punjabi for to fry, then the Punjabi for to roast.
\end{quote}

\subsection*{You'll want to know: \pa{ਕੜਛੀ}}
\textbf{\pa{ਕੜਛੀ}} — \emph{kaṛchhī} — \textquotedblleft{}a ladle\textquotedblright{}.

\begin{quote}
\pa{ਕੜਛੀ} \pa{ਪਤੀਲੇ} \pa{ਵਿੱਚ} \pa{ਹੈ।}
\end{quote}

\begin{quote}
\emph{kaṛchhī patīle vich hai.}
\end{quote}

\begin{quote}
The ladle is in the cooking pot.
\end{quote}

A feminine noun.

\begin{cousinweb}
\textbf{\pa{ਕੜਛੀ}} (\emph{kaṛchhī}) is inherited from a reconstructed \emph{*kaṭacchu}, \textquotedblleft{}ladle\textquotedblright{}; CDIAL lists Punjabi \emph{kaṛchī}, \textquotedblleft{}iron ladle\textquotedblright{} (CDIAL 2633).
\end{cousinweb}

\begin{grammarlens}[title={What you've built}]
One more word for this chapter.
\end{grammarlens}

\subsection*{Your turn}
\begin{itemize}
\raggedright
  \item \emph{Say it:} \emph{kaṛchhī}
  \item \emph{Say it:} \emph{kaṛchhī}, once more
  \item \emph{Say it:} \emph{kaṛchhī patīle vich hai.} — The ladle is in the cooking pot.
  \item \emph{Recall:} say the Punjabi for to fry, then the Punjabi for to roast, then say \emph{kaṛchhī} again
  \item \emph{Write it:} copy \textbf{\pa{ਕੜਛੀ}} once, checking the whole word
\end{itemize}

\begin{tcolorbox}[breakable,colback=teal!4,colframe=teal!35!black,title={Before you move on}]
What does \pa{ਕੜਛੀ} mean? (\textquotedblleft{}A ladle\textquotedblright{}.) Say it once more.
\end{tcolorbox}
\section[\texorpdfstring{\pa{ਛਿੜਕਣਾ} (chhiṛakṇā)}{ਛਿੜਕਣਾ (chhiṛakṇā)}]{\pa{ਛਿੜਕਣਾ} (chhiṛakṇā) — to sprinkle}

\label{lesson:PA-C242-chhirakna}



\begin{quote}
Before the new one: say the Punjabi for to roast, then the Punjabi for a ladle.
\end{quote}

\subsection*{You'll want to know: \pa{ਛਿੜਕਣਾ}}
\textbf{\pa{ਛਿੜਕਣਾ}} — \emph{chhiṛakṇā} — \textquotedblleft{}to sprinkle\textquotedblright{}.

\begin{quote}
\pa{ਮੈਂ} \pa{ਸਬਜ਼ੀ} \pa{ਉੱਤੇ} \pa{ਲੂਣ} \pa{ਛਿੜਕਦੀ} \pa{ਹਾਂ।}
\end{quote}

\begin{quote}
\emph{maiṁ sabzī utte lūṇ chhiṛakdī hāṁ.}
\end{quote}

\begin{quote}
I sprinkle salt on the vegetables.
\end{quote}

Past (I …): \pa{ਮੈਂ} \pa{ਲੂਣ} \pa{ਛਿੜਕਿਆ}, \emph{maiṁ lūṇ chhiṛkiā}.

\begin{cousinweb}
\textbf{\pa{ਛਿੜਕਣਾ}} (\emph{chhiṛakṇā}) is inherited from a reconstructed \emph{*chiṭ}, \textquotedblleft{}sudden movement, splash, scatter\textquotedblright{}; CDIAL lists Punjabi \emph{chiṛkaṇā}, \textquotedblleft{}to sprinkle\textquotedblright{} (CDIAL 5035).
\end{cousinweb}

\begin{grammarlens}[title={What you've built}]
One more word for this chapter.
\end{grammarlens}

\subsection*{Your turn}
\begin{itemize}
\raggedright
  \item \emph{Say it:} \emph{chhiṛakṇā}
  \item \emph{Say it:} \emph{chhiṛakṇā}, once more
  \item \emph{Say it:} \emph{maiṁ sabzī utte lūṇ chhiṛakdī hāṁ.} — I sprinkle salt on the vegetables.
  \item \emph{Recall:} say the Punjabi for to roast, then the Punjabi for a ladle, then say \emph{chhiṛakṇā} again
  \item \emph{Write it:} copy \textbf{\pa{ਛਿੜਕਣਾ}} once, checking the whole word
\end{itemize}

\begin{tcolorbox}[breakable,colback=teal!4,colframe=teal!35!black,title={Before you move on}]
What does \pa{ਛਿੜਕਣਾ} mean? (\textquotedblleft{}To sprinkle\textquotedblright{}.) Say it once more.
\end{tcolorbox}
\section[\texorpdfstring{\pa{ਸੜਨਾ} (saṛnā)}{ਸੜਨਾ (saṛnā)}]{\pa{ਸੜਨਾ} (saṛnā) — to get burnt}

\label{lesson:PA-C242-sarna}



\begin{quote}
Before the new one: say the Punjabi for a ladle, then the Punjabi for to sprinkle.
\end{quote}

\subsection*{You'll want to know: \pa{ਸੜਨਾ}}
\textbf{\pa{ਸੜਨਾ}} — \emph{saṛnā} — \textquotedblleft{}to get burnt\textquotedblright{}.

\begin{quote}
\pa{ਜੇ} \pa{ਤਵਾ} \pa{ਬਹੁਤ} \pa{ਗਰਮ} \pa{ਹੈ} \pa{ਤਾਂ} \pa{ਰੋਟੀ} \pa{ਸੜਦੀ} \pa{ਹੈ।}
\end{quote}

\begin{quote}
\emph{je tavā bahut garam hai tāṁ roṭī saṛdī hai.}
\end{quote}

\begin{quote}
If the griddle is too hot, the roti burns.
\end{quote}

Past (it …): \pa{ਉਹ} \pa{ਰੋਟੀ} \pa{ਸੜੀ}, \emph{oh roṭī saṛī}.

\begin{cousinweb}
\textbf{\pa{ਸੜਨਾ}} (\emph{saṛnā}) is inherited from Sanskrit \emph{śaṭati}, \textquotedblleft{}disintegrates\textquotedblright{}; CDIAL lists Punjabi \emph{saṛnā}, \textquotedblleft{}to rot, be burnt\textquotedblright{} (CDIAL 12268).
\end{cousinweb}

\begin{grammarlens}[title={What you've built}]
One more word for this chapter.
\end{grammarlens}

\subsection*{Your turn}
\begin{itemize}
\raggedright
  \item \emph{Say it:} \emph{saṛnā}
  \item \emph{Say it:} \emph{saṛnā}, once more
  \item \emph{Say it:} \emph{je tavā bahut garam hai tāṁ roṭī saṛdī hai.} — If the griddle is too hot, the roti burns.
  \item \emph{Recall:} say the Punjabi for a ladle, then the Punjabi for to sprinkle, then say \emph{saṛnā} again
  \item \emph{Write it:} copy \textbf{\pa{ਸੜਨਾ}} once, checking the whole word
\end{itemize}

\begin{tcolorbox}[breakable,colback=teal!4,colframe=teal!35!black,title={Before you move on}]
What does \pa{ਸੜਨਾ} mean? (\textquotedblleft{}To get burnt\textquotedblright{}.) Say it once more.
\end{tcolorbox}
